\documentclass[11pt,a4paper]{ctexart}

% ============================ 字体与页面 ============================
\usepackage{fontspec}
\setmainfont{Arial}
\setmonofont{Consolas}
\setCJKmainfont{Microsoft YaHei}[AutoFakeBold=3.0]
\setCJKsansfont{Microsoft YaHei}
\setCJKmonofont{Noto Sans SC}

\usepackage[margin=2.1cm]{geometry}
\linespread{1.35}
\emergencystretch=3em
\setlength{\parindent}{2em}
\setlength{\parskip}{0.35em}

% ============================ 颜色 ============================
\usepackage{xcolor}
\definecolor{hlblue}{RGB}{31,78,121}     % 深蓝（主标题/规则）
\definecolor{accent}{RGB}{0,112,192}     % 亮蓝（小节）
\definecolor{tblhead}{RGB}{221,235,247}  % 表头底色（浅蓝）
\definecolor{tblrule}{RGB}{120,150,190}  % 表格规则线
\definecolor{pgreen}{RGB}{46,125,50}
\definecolor{porange}{RGB}{230,81,0}
\definecolor{pred}{RGB}{198,40,40}
\definecolor{pgray}{RGB}{120,120,120}

% ============================ 标题层级 ============================
\usepackage{titlesec}
\titleformat{\section}
  {\Large\bfseries\color{hlblue}}{\thesection.}{0.6em}{}
  [\vspace{-0.35em}{\color{hlblue}\titlerule[1.3pt]}]
\titlespacing*{\section}{0pt}{1.4em}{0.7em}

\titleformat{\subsection}
  {\large\bfseries\color{accent}}{\thesubsection}{0.6em}{}
\titlespacing*{\subsection}{0pt}{1.1em}{0.45em}

\titleformat{\subsubsection}
  {\normalsize\bfseries\color{hlblue!85!black}}{\thesubsubsection}{0.6em}{}
\titlespacing*{\subsubsection}{0pt}{0.9em}{0.35em}

% ============================ 页眉页脚 ============================
\usepackage{fancyhdr}
\usepackage{lastpage}
\setlength{\headheight}{15pt}
\pagestyle{fancy}
\fancyhf{}
\fancyhead[L]{\small\color{pgray} RR-GID 项目阶段性汇报}
\fancyhead[R]{\small\color{pgray} 2026-08-26}
\fancyfoot[L]{\small\color{pgray} 基于唯一规范 RR\_GID\_CN.pdf}
\fancyfoot[C]{\small 第 \thepage\ 页\ /\ 共 \pageref{LastPage}\ 页}
\renewcommand{\headrulewidth}{0.6pt}
\renewcommand{\headrule}{\color{hlblue}\hrule width\headwidth height 0.6pt}

% ============================ 表格 ============================
\usepackage{booktabs}
\usepackage{longtable}
\usepackage{array}
\usepackage{colortbl}
\newcolumntype{C}[1]{>{\centering\arraybackslash}p{#1}}
\newcolumntype{L}[1]{>{\raggedright\arraybackslash}p{#1}}
\arrayrulecolor{tblrule}
\renewcommand{\arraystretch}{1.45}
\setlength{\tabcolsep}{7pt}
\setlength{\LTpre}{0.8em}
\setlength{\LTpost}{0.8em}
\newcommand{\hcell}[1]{\textcolor{white}{\bfseries #1}}

% ============================ 高亮框 ============================
\usepackage[most]{tcolorbox}
\newtcolorbox{keybox}[2]{%
  breakable, arc=2mm, boxrule=0.9pt,
  left=9pt, right=9pt, top=7pt, bottom=7pt,
  colback=#2!6, colframe=#2, coltitle=white, colbacktitle=#2,
  fonttitle=\bfseries\small, title=#1, before skip=0.8em, after skip=0.8em}

% ============================ 列表与数学 ============================
\usepackage{amsmath,amssymb}
\usepackage{enumitem}
\setlist{topsep=2pt, itemsep=2pt, parsep=0pt}
\usepackage[colorlinks=true, linkcolor=hlblue, urlcolor=accent, citecolor=hlblue]{hyperref}

% ============================ 文档 ============================
\begin{document}

% ---------- 标题块 ----------
\begin{center}
  {\color{pgray}\small RR-GID\_CN · Reference-Relative Generative Information Design}
  \par\vspace{0.8em}
  {\color{hlblue}\rule{0.94\textwidth}{1.4pt}}
  \par\vspace{1.1em}
  {\LARGE\bfseries\color{hlblue} RR-GID 项目进展汇报}\par\vspace{0.5em}
  {\large\bfseries 已完成工作与 P4 理论阻塞分析}\par\vspace{0.9em}
  {\color{pgray}\small 阶段性汇报 · 2026-08-26 · 基于项目唯一规范 \texttt{RR\_GID\_CN.pdf}}\par\vspace{0.6em}
  {\color{hlblue}\rule{0.94\textwidth}{1.4pt}}
\end{center}
\vspace{0.6em}

% ================================================================
\section{PDF 任务分成了哪些步骤}
% ================================================================
PDF 是一个“参考分布驱动的生成式主动设计”方法（RR-GID）。工程上按 \texttt{docs/PROJECT\_PLAN.md} 拆成 P0–P12 十二个阶段，逻辑上是四段：

\begin{longtable}{@{}C{2.1cm}C{1.9cm}L{8.6cm}C{2.4cm}@{}}
  \caption{PDF 任务的 P0–P12 阶段划分}\label{tbl:stages}\\
  \toprule[1.4pt]
  \rowcolor{tblhead}\hcell{段落} & \hcell{阶段} & \hcell{内容} & \hcell{性质} \\
  \midrule[0.8pt]
  \endfirsthead
  \toprule[1.4pt]
  \rowcolor{tblhead}\hcell{段落} & \hcell{阶段} & \hcell{内容} & \hcell{性质} \\
  \midrule[0.8pt]
  \endhead
  \bottomrule[1.4pt]
  \endfoot
  工程地基 & P0--P3 & 环境/Git、Synthetic 精确 oracle（16 维 GMM + sinh warp、120 panels、12 维 feature map）、四种 allocation policy、RR-GID 的 pilot--design--update 完整算法 & 纯工程 \\
   Synthetic 主实验 & P4--P7 & P4 S1（三策略最优性门禁）、P5 最终四策略、P6 冻结 VAEAC 生成器、P7 S2（非线性 + 生成器复用） & \textbf{PDF 核心理论验证} \\
  真实数据 & P8--P10 & P8 Gas 数据预处理、P9 R1 半合成、P10 R2 自然漂移 & 真实数据验证 \\
  收尾 & P11--P12 & 统一图表、最终复现与冻结 & 产出 \\
  \bottomrule[1.4pt]
\end{longtable}

其中 \textbf{P4 是整个 PDF 的理论门禁}：它要求 Synthetic 上三策略的 $B\cdot\mathrm{KL}$ 收敛到理论常数、design ratio 接近 1。\textbf{P4 不过，P5--P12 都不能进入}。

% ================================================================
\section{每个步骤完成了什么}
% ================================================================

\subsection{已完成并验证（P0--P3 工程 + 算法）}
\begin{itemize}
  \item \textbf{Synthetic 精确 oracle}：16 维四成分 GMM + sinh warp（$\alpha=1$）、120 个坐标对 panels、12 维 bounded feature map、精确 Gaussian-mixture conditional sampler。
  \item \textbf{三种 baseline policy}：Uniform SQD、A-OSQD（完整参考协方差 A-optimal）、oracle RR-GID。
  \item \textbf{RR-GID 完整算法}：balanced pilot（$b_B=\lceil 10B^{1/3}\rceil$）、HT moment、cross-completion 信息估计 + PSD 投影、cost-aware Frank-Wolfe design、$J=2$ 步 Fisher-scoring。
  \item \textbf{预算守恒}、配对 target draws、seed 方案、KL/design ratio 计算。
  \item 单元测试在当前环境 \textbf{32 passed}。
\end{itemize}

\subsection{有工程资产但未通过正式验收（P5--P10）}
\begin{itemize}
  \item \textbf{P5/P7/P9/P10}：旧结果文件数量正确（如 P10 三 campaign $\times$ 4 budgets $\times$ 50 reps $\times$ 4 policies = 2400 行），但旧 acceptance 明确记为 \textcolor{pred}{\textbf{NOT PASSED}}，且 P9 final estimator 曾用 empirical kernel 而非冻结 VAEAC、P10 指标（Bregman projection loss、policy-specific $\hat{\beta}$、32 held-out functions）未按 PDF Eq.17 实现。这些\textbf{不能作为最终结论}。
  \item \textbf{P6}：VAEAC 生成器经 codex 11 次重训，无条件矩仍未达质量门禁（无条件 std 误差 $\sim$10--17、条件 RMSE $\sim$6），generator 质量不足以支撑依赖它的 P5/P7/P9。
\end{itemize}

\subsection{未完成（P11--P12）}
依赖 P4--P10 通过，尚未进入。

% ================================================================
\section{P4 卡在哪：理论上的问题}
% ================================================================

\subsection{想要的结果（PDF 理论预期）}
PDF Sec.5 的局部渐近最优性要求，对固定 target $\beta^*$、oracle design $p^*$：
\begin{equation}
  \mathbb{E}\big[ \mathrm{KL}\big( Q_{\beta^*}\,\lVert\, Q_{\hat{\beta}_B} \big) \big]
  = \frac{1}{2B}\,\mathrm{tr}\!\big( F(\beta^*)\, M(p^*)^{-1} \big) + o(1/B)
  \label{eq:kl}
\end{equation}
即
\[
  B\cdot\mathrm{KL} \;\longrightarrow\; \tfrac12\,\Phi(p^*)
  \qquad\text{（一个与 $B$ 无关的理论常数，如 $B=2000$ 时 $\tfrac12\Phi\approx 31.9$）}
\]
且
\[
  \text{design ratio} \;=\; \frac{\Phi(\hat p)}{\Phi(p^*)} \;\longrightarrow\; 1 .
\]

\subsection{实际得到的结果（formal replication）}
codex 在新开发机（8$\times$A800）上完成的部分正式预算（每档 200 reps $\times$ 3 policies，paired draws）：

\begin{longtable}{@{}C{1.6cm}C{4.2cm}C{2.6cm}C{2.6cm}C{2.6cm}@{}}
  \caption{P4 formal replication：各预算下三种策略的 $B\cdot\mathrm{KL}$}\label{tbl:p4formal}\\
  \toprule[1.4pt]
  \rowcolor{tblhead}\hcell{Budget} & \hcell{oracle RR-GID 均值} & \hcell{Uniform} & \hcell{A-OSQD} & \hcell{理论线 $\tfrac12\Phi$} \\
  \midrule[0.8pt]
  \endfirsthead
  \toprule[1.4pt]
  \rowcolor{tblhead}\hcell{Budget} & \hcell{oracle RR-GID 均值} & \hcell{Uniform} & \hcell{A-OSQD} & \hcell{理论线 $\tfrac12\Phi$} \\
  \midrule[0.8pt]
  \endhead
  \bottomrule[1.4pt]
  \endfoot
  2000 & 44.11 (CI [40.8, 47.4]) & 82.90 & 128.33 & \textbf{31.9} \\
  4000 & 58.35 (CI [53.2, 63.5]) & 96.49 & 180.91 & \textbf{63.8} \\
  \bottomrule[1.4pt]
\end{longtable}

（另一组早期未用正式配置的结果，$B\cdot\mathrm{KL}$ 随预算线性上升到 1795.7 @ $B=32000$，design ratio 高达 42.7；codex 已判定该组 estimator 实现不符 PDF Alg.2，不作为正式结论。）

\begin{keybox}{核心矛盾}{pred}
\begin{itemize}
  \item oracle RR-GID 的 $B\cdot\mathrm{KL}$ \textbf{没有随预算收敛到理论常数}，而是停留在 $\sim$44--58 的平台附近（$B=2000$ 略高于理论线 31.9，$B=4000$ 反而低于理论线 63.8）。
  \item \textbf{KL 没有按 $O(1/B)$ 衰减}（KL 停在 $\sim$0.02--0.03 量级），而是 $B\cdot\mathrm{KL}\approx$ 常数。这是 $B\cdot\mathrm{KL}$ 不收敛的直接表现。
  \item design ratio 明显 $>1$。
\end{itemize}
\end{keybox}

\subsection{为什么卡住：理论假设与 PDF 冻结 pilot 公式的张力}
通过固定 seed 的小规模诊断，codex 已经把问题\textbf{收窄到很具体的一环}：
\begin{enumerate}
  \item \textbf{$\beta$ 求解器本身是精确的}：用精确的 target moment $\mu_{\beta^*}$ 反解 $\beta$，误差 $\approx 6.3\times10^{-9}$。问题不在求解器。
  \item \textbf{卡点在 balanced HT pilot}：PDF 强制 pilot 预算 $b_B=\lceil 10B^{1/3}\rceil$（$B=8000$ 时只有 $\sim$200 观测），要在 120 个 panels 上估计 12 维 $\beta$。在\textbf{弱 Fisher 方向}（某些 feature 变化对 KL 影响很小），pilot 观测覆盖不足，HT moment 产生\textbf{不随 $B$ 衰减的系统偏差}（单 seed $\beta$ 误差约 1.4--2.0）。
  \item \textbf{这个偏差是 $O(1)$，不是理论要求的 $O(B^{-1/6})$}。PDF 理论假设 $\lVert\hat\beta_p-\beta^*\rVert = O_p(B^{-\gamma/2})$（$\gamma=1/3$），但实际 pilot 的系统偏差不随 $B$ 缩小。
  \item \textbf{偏差被后续 Fisher-scoring 放大}：$\hat\beta$ 有偏 $\Rightarrow$ 估计的信息矩阵 $H$ 和观测得分 $U$ 都带噪声 $\Rightarrow$ Newton 步在弱方向过冲。
\end{enumerate}

\subsection{数学依据（为什么这不是数值 bug）}
设 pilot 的系统偏差为常数 $\delta=\hat\beta_p-\beta^*$（不随 $B$ 衰减）。则最终估计的 KL 可展开为：
\[
  \mathrm{KL}\big( Q_{\beta^*}\,\lVert\, Q_{\hat\beta} \big)
  \;\approx\; \tfrac12\,(\hat\beta-\beta^*)^\top\, F(\beta^*)\,(\hat\beta-\beta^*)
  \qquad\text{（二阶 Taylor，$F$ 为 Fisher 信息）}
\]
若 $\hat\beta$ 的偏差以常数 $\delta$ 为主导（不随 $B$ 收敛到 0），则：
\[
  B\cdot\mathrm{KL}
  \;\approx\; B\cdot\tfrac12\,\delta^\top F\,\delta
  \;=\; \tfrac12\,\Phi(p^*) \;+\; \frac{B\cdot(\delta^\top F\,\delta)}{2}
\]
即 \textbf{$B\cdot\mathrm{KL}$ 中混入了一个随 $B$ 线性增长的项 $B\cdot(\delta^\top F\,\delta)/2$}。只要 $\delta\neq 0$（常数级偏差），$B\cdot\mathrm{KL}$ 就不可能收敛到理论常数。这正是观测到的现象：KL 停在常数量级 $\Rightarrow$ $B\cdot\mathrm{KL}$ 随 $B$ 上升。

\begin{keybox}{关键推论}{hlblue}
这不是统计波动（加大 replication 不会消失），也不是 MC 尺寸不足（加大 $L_U/H_{\text{tilted}}/H_{\text{cond}}$ 只降低估计方差，不消除 pilot 系统偏差）。它是 \textbf{PDF 冻结的 pilot 预算公式与理论相合性假设之间的有限样本张力}。
\end{keybox}

% ================================================================
\section{为解决问题做的所有尝试（按调整思路分组）}
% ================================================================
以下按“想解决什么问题、怎么做、结果如何”组织。全部保留了日志/诊断结果，未覆盖任何正式结果。

\subsection{尝试修复 pilot 偏差（核心方向）}
\begin{longtable}{@{}L{4.2cm}L{4.6cm}L{6.2cm}@{}}
  \caption{修复 pilot 偏差的尝试}\label{tbl:pilot}\\
  \toprule[1.4pt]
  \rowcolor{tblhead}\hcell{尝试} & \hcell{思路} & \hcell{结果} \\
  \midrule[0.8pt]
  \endfirsthead
  \toprule[1.4pt]
  \rowcolor{tblhead}\hcell{尝试} & \hcell{思路} & \hcell{结果} \\
  \midrule[0.8pt]
  \endhead
  \bottomrule[1.4pt]
  \endfoot
  pilot moment 用 reference Fisher 稳定初始化 & 减少 pilot 弱方向的初始化误差 & 单 seed 略改善，不收敛 \\
  pilot $\beta$ 投影到 compact Theta & 约束 pilot 解不越界 & $\hat\beta_{\text{norm}}$ 从发散到 2.19 稳定，但 KL 仍不衰减 \\
  加大诊断 pilot 预算（非冻结配置） & 验证是否 pilot 预算过小 & 部分 seed 改善，但不是 PDF 冻结实验 \\
  恢复历史 \texttt{use\_oracle\_H=true} & 绕开估计 $H$ 的噪声 & \textbf{反而更差}，未采用 \\
  \bottomrule[1.4pt]
\end{longtable}

\subsection{尝试修复 Fisher-scoring 步长稳定性}
\begin{longtable}{@{}L{4.2cm}L{4.6cm}L{6.2cm}@{}}
  \caption{Fisher-scoring 步长稳定性的尝试}\label{tbl:scoring}\\
  \toprule[1.4pt]
  \rowcolor{tblhead}\hcell{尝试} & \hcell{思路} & \hcell{结果} \\
  \midrule[0.8pt]
  \endfirsthead
  \toprule[1.4pt]
  \rowcolor{tblhead}\hcell{尝试} & \hcell{思路} & \hcell{结果} \\
  \midrule[0.8pt]
  \endhead
  \bottomrule[1.4pt]
  \endfoot
  固定步长 0.25/0.5/0.75/1.0 & 抑制 Newton 过冲 & 无一致胜者，不构成修复 \\
  trust-region & 限制更新范数 & 引入持续偏差，反而恶化 \\
  恢复 PDF full step + Theta/L1 投影 & 回到 Alg.2 原定义 & 弱方向仍过冲 \\
  \bottomrule[1.4pt]
\end{longtable}

\subsection{尝试降低信息估计噪声（H 层）}
\begin{longtable}{@{}L{4.2cm}L{4.6cm}L{6.2cm}@{}}
  \caption{信息估计噪声（H 层）的尝试}\label{tbl:info}\\
  \toprule[1.4pt]
  \rowcolor{tblhead}\hcell{尝试} & \hcell{思路} & \hcell{结果} \\
  \midrule[0.8pt]
  \endfirsthead
  \toprule[1.4pt]
  \rowcolor{tblhead}\hcell{尝试} & \hcell{思路} & \hcell{结果} \\
  \midrule[0.8pt]
  \endhead
  \bottomrule[1.4pt]
  \endfoot
  cross-completion + PSD 投影 & 消除自归一化偏差 & 保留（正确实现），但不消除 pilot 偏差 \\
  Sobol QMC 条件积分（\texttt{b75720a} 系列 6 个 commit） & 降低条件期望 MC 方差 & 方差改善，偏差不消 \\
  information MC budget-growing（$L_U/B\to\infty$） & 满足 PDF Alg.2 的 $L_{U,B}\to\infty$ & $L_U$ 效果非单调、小于 seed 方差；$H$ 增长有微弱信号但不稳定 \\
  缓存 conditional Gaussian Cholesky + feature-mean batch & 纯性能优化 & 快很多，不影响统计 \\
  \bottomrule[1.4pt]
\end{longtable}

\subsection{性能优化（消除工程瓶颈，非统计）}
\begin{itemize}
  \item exact conditional rejection 逐观测调用 $\rightarrow$ 向量化 + Cholesky 缓存（单 rep 从数分钟降到 $\sim$1.5 分钟）
  \item 并行 shard 重复计算 prepared oracle $\rightarrow$ 显式复用已审计 pkl
  \item 设备绑定 \texttt{RR\_GID\_CN\_CUDA\_DEVICE}（8 卡各自绑定，GPU0/1 同 seed 输出完全一致）
\end{itemize}

\subsection{P6 VAEAC 质量（11 次重训，依赖它的阶段被阻塞）}
\begin{longtable}{@{}L{7.2cm}L{7.4cm}@{}}
  \caption{P6 VAEAC 生成器质量的尝试}\label{tbl:p6}\\
  \toprule[1.4pt]
  \rowcolor{tblhead}\hcell{尝试} & \hcell{结果} \\
  \midrule[0.8pt]
  \endfirsthead
  \toprule[1.4pt]
  \rowcolor{tblhead}\hcell{尝试} & \hcell{结果} \\
  \midrule[0.8pt]
  \endhead
  \bottomrule[1.4pt]
  \endfoot
  四成分 GMM prior + KMeans 初始化 + prior 冻结 & 无条件 std 误差从 57 $\rightarrow$ 10 \\
  增强 panel-conditioned mask 训练 + 提高条件权重 & 条件 RMSE 仍 $\sim$6 \\
  400/800 epoch 重训、\texttt{z\_std} 归一化修正 & 无条件均值误差降到 0.95，std 误差 10，条件 RMSE 6 \\
  \bottomrule[1.4pt]
\end{longtable}

P6 生成器仍未达质量门禁，是除 P4 外的第二个阻塞项。

% ================================================================
\section{想要的结果和现在的差距}
% ================================================================
\begin{longtable}{@{}L{3.4cm}L{4.2cm}L{5.0cm}C{2.6cm}@{}}
  \caption{理论目标与实际结果的差距}\label{tbl:gap}\\
  \toprule[1.4pt]
  \rowcolor{tblhead}\hcell{维度} & \hcell{理论想要} & \hcell{现在} & \hcell{差距} \\
  \midrule[0.8pt]
  \endfirsthead
  \toprule[1.4pt]
  \rowcolor{tblhead}\hcell{维度} & \hcell{理论想要} & \hcell{现在} & \hcell{差距} \\
  \midrule[0.8pt]
  \endhead
  \bottomrule[1.4pt]
  \endfoot
  $B\cdot\mathrm{KL}$ 收敛 & $\to \tfrac12\Phi(p^*)$（常数平台） & 停在 $\sim$44--58 平台，不随 $B$ 衰减 & \textbf{不收敛} \\
  KL 衰减率 & $O(1/B)$ & 停在 $\sim$0.02--0.03 & \textbf{不衰减} \\
  design ratio & $\to 1$ & $>1$ & \textbf{偏高} \\
  oracle RR-GID 相对最优 & 稳定优于 Uniform/A-OSQD & 部分预算不显著 & \textbf{不稳定} \\
  \bottomrule[1.4pt]
\end{longtable}

% ================================================================
\section{现有实验说明 PDF 理论可能哪里需要审视（数学依据）}
% ================================================================
以下是从实验结果\textbf{反推} PDF 理论里可能需要更仔细审视的假设，不是否定 PDF。

\subsubsection{1. pilot 相合性假设可能过强}
PDF Sec.5.3 假设 $\lVert\hat\beta_p-\beta^*\rVert = O_p(B^{-\gamma/2})$（$\gamma=1/3$），从而 $\Phi(\hat p)\to \min_p \Phi(p)$。但 PDF 同时冻结 $b_B=\lceil 10B^{1/3}\rceil$——在这个 pilot 预算下，HT moment 在 120 个高维弱信息面板上的系统偏差是 \textbf{$O(1)$ 而非 $O(B^{-1/6})$}（实测单 seed $\beta$ 误差 1.4--2.0，不随 $B$ 衰减）。数学上，若 $\delta=O(1)$，则 $B\cdot\mathrm{KL}=\tfrac12\Phi(p^*)+B\cdot(\delta^\top F\delta)/2+o(B)$，收敛被线性项破坏。

\textbf{可能的修订方向}（供讨论，不擅自改）：
\begin{itemize}
  \item 增大 pilot 预算（如 $b_B=\lceil B^{1/2}\rceil$ 或与 feature 维度/panel 数挂钩）；
  \item 或 pilot 用 reference-Fisher 定向分配（优先覆盖强信息方向），而非均匀 balanced；
  \item 或在理论部分补充对 pilot 有限样本偏差的修正项分析。
\end{itemize}

\subsubsection{2. 有限样本信息矩阵估计的 operator error 条件}
PDF Alg.2 要求 update 阶段的信息估计 $\hat I_S$ 的 operator error 消失，但固定 MC 尺寸（$h_{\text{tilted}}, h_{\text{cond}}$）不满足渐近条件。codex 的 H-growth 诊断（$(256,8)\to(1024,32)$）显示微弱改善但 CI 太宽，\textbf{尚未证明误差消失}。

\subsubsection{3. KL 的 O(1/B) 率依赖 pilot 一致性}
由于偏差项不消失，KL 无法达到 PDF Eq.13 的 $\mathbb{E}[\mathrm{KL}]=O(1/B)$。这是\textbf{有限样本下的理论预期偏差}，不是实现错误。

% ================================================================
\section{当前状态与下一步}
% ================================================================
\begin{keybox}{诚实结论}{pgreen}
项目在工程上完成了 P0--P3 和大部分实验代码/结果框架；但 \textbf{P4 的理论门禁未通过，P5--P12 不能进入}。P4 卡在“PDF 冻结的 pilot 预算公式与理论相合性假设的有限样本张力”，codex 已穷尽 PDF 框架内的修法（pilot 初始化、oracle-H、阻尼、trust-region、QMC、H-growth）均未收敛。P6 生成器质量是第二个阻塞项。
\end{keybox}

\textbf{下一步选项}（可能的决策）：
\begin{enumerate}
  \item \textbf{讨论是否可修订 PDF 的 pilot 预算}（最直接恢复收敛）；
  \item \textbf{把有限样本 pilot 偏差作为现象写进论文}（给出 $B\cdot\mathrm{KL}=\text{常数}+B\cdot\text{bias}^2/2$ 的精确形式 + 偏差修正），不强求收敛到理论常数；
  \item \textbf{调整论文重心}到不依赖 P4 收敛的相对结论（RR-GID 相对 Uniform/A-OSQD 的 compute 优势、reuse frontier）。
\end{enumerate}

{\sloppy 代码、结果、诊断与交接文档已全部推送到 GitHub：\url{https://github.com/15628925702/RR_GID_CN.git}（main 分支）。交接文档：\url{docs/PERSISTENT_CONTEXT_20260826.md}、\url{docs/HANDOVER_STATUS_dev.md}。\par}

\end{document}
